\begin{problem}[The Calkin--Wilf Tree and Unimodular Matrix Walks]
% \textit{Target level: Senior Specialist Mathematics / Olympiad bridge.}

Let $(x_n)_{n\geq 1}$ be defined by
\[
x_1=1,\qquad x_{n+1}=\frac{1}{2\lfloor x_n\rfloor-x_n+1}.
\]
Write $x_n=p_n/q_n$ in lowest terms, where $p_n,q_n\in\mathbb Z^+$.

\begin{enumerate}[label=\textbf{(\alph*)}]
\item Compute $x_2,\ldots,x_7$ and arrange them in the first three levels of a binary tree. Prove, by induction, that
\[
x_{2k}=\frac{x_k}{x_k+1},\qquad x_{2k+1}=x_k+1.
\]

\item Define
\[
M_0=\begin{pmatrix}1&0\\1&1\end{pmatrix},\qquad
M_1=\begin{pmatrix}1&1\\0&1\end{pmatrix}.
\]
Show that, with $\mathbf v_k=(p_k,q_k)^T$,
\[
\mathbf v_{2k+b}=M_b\mathbf v_k\qquad(b\in\{0,1\}).
\]
Hence, if $n=(1b_{m-1}\cdots b_0)_2$, prove that
\[
\mathbf v_n=M_{b_0}M_{b_1}\cdots M_{b_{m-1}}\binom11.
\]
Use the inverse operations to prove that every positive rational number occurs exactly once in the sequence.

\item Put
\[
u_k=\frac1{p_k^2q_k}+\frac1{p_kq_k^2},\qquad
S_n=\sum_{k=1}^n\frac1{p_kq_k(p_k+q_k)}.
\]
Prove the tree-telescoping identity
\[
\frac1{p_kq_k(p_k+q_k)}=u_k-u_{2k}-u_{2k+1},
\]
and deduce that
\[
S_n=2-\sum_{j=n+1}^{2n+1}u_j.
\]
Conclude that $(S_n)$ is strictly increasing and bounded above by $2$.

\item Let
$\mathcal C=\{(p,q)\in\mathbb Z^+\times\mathbb Z^+:\gcd(p,q)=1\},\quad
T_N=\{(p,q)\in\mathcal C:p+q\leq N\},
$
and define $\Sigma_N=\sum_{(p,q)\in T_N}1/[pq(p+q)]$. Show that
\[
0<2-\Sigma_N<\frac8N\sum_{m=1}^{N-1}\frac1m.
\]
Explain why this proves the exact value of $\lim_{n\to\infty}S_n$.
\end{enumerate}
\end{problem}

\begin{hint}
\begin{itemize}
\item Use $2\lfloor x_k\rfloor-x_k+1=1/x_{k+1}$ in the induction. The even child lies in $(0,1)$.
\item $M_0^{-1}(p,q)^T=(p,q-p)^T$ when $p<q$, while $M_1^{-1}(p,q)^T=(p-q,q)^T$ when $p>q$. Each subtraction lowers $p+q$.
\item Substitute $(p_{2k},q_{2k})=(p_k,p_k+q_k)$ and $(p_{2k+1},q_{2k+1})=(p_k+q_k,q_k)$, then sum over $k$.
\item For the boundary estimate, set $m=\min(a,b)$ and $M=\max(a,b)$. A boundary point satisfies $N-m<M\leq N$, so there are at most $2m$ ordered pairs for each $m$.
\end{itemize}
\end{hint}

\begin{solution}
\textbf{(a)} The first terms are
\[
1,\ \frac12,\ 2,\ \frac13,\ \frac32,\ \frac23,\ 3.
\]
Thus the children of $x_k$ are $x_k/(x_k+1)$ and $x_k+1$. The assertion holds for $k=1$. If it holds for $k$, then
\[
x_{2k+2}=\frac1{2(\lfloor x_k\rfloor+1)-(x_k+1)+1}
=\frac{x_{k+1}}{x_{k+1}+1},
\]
which lies in $(0,1)$, and therefore $x_{2k+3}=x_{k+1}+1$. This completes the induction.

\textbf{(b)} The child formulae give
\[
(p_{2k},q_{2k})=(p_k,p_k+q_k),\qquad
(p_{2k+1},q_{2k+1})=(p_k+q_k,q_k),
\]
and no cancellation occurs because the parent coordinates are coprime. These are exactly the two matrix equations. Repeatedly removing the least significant binary digit gives the displayed product.

For a coprime pair $(p,q)\ne(1,1)$, exactly one of $p<q$ and $p>q$ holds. Apply the corresponding inverse subtraction. Positivity and coprimality are preserved, while $p+q$ decreases, so the process terminates at $(1,1)$ and is unique. Hence the sequence enumerates $\mathbb Q^+$ bijectively.

\textbf{(c)} Direct substitution of the two child pairs gives
\[
u_k-u_{2k}-u_{2k+1}=\frac1{p_kq_k(p_k+q_k)}.
\]
Summing for $1\leq k\leq n$ cancels the indices $2,\ldots,2n+1$ and yields
\[
S_n=u_1-\sum_{j=n+1}^{2n+1}u_j=2-\sum_{j=n+1}^{2n+1}u_j<2.
\]
The summands defining $S_n$ are positive, so $S_n$ is strictly increasing.

\textbf{(d)} The bijection in (b) and positivity allow the series to be re-indexed by $p+q\leq N$; hence $\lim S_n=\lim\Sigma_N$. Tree telescoping over $T_N$ leaves the root contribution $u(1,1)=2$ minus the children outside $T_N$ whose parent is inside. These boundary children are precisely the pairs $(a,b)$ with $a+b>N$ and $\max(a,b)\leq N$.

Writing $m=\min(a,b)$ and $M=\max(a,b)$, there are at most $2m$ ordered pairs for each $m$, and $M>N/2$. Therefore each contributes less than $4/(m^2N)$, giving
\[
0<2-\Sigma_N<\sum_{m=1}^{N-1}2m\frac4{m^2N}
=\frac8N\sum_{m=1}^{N-1}\frac1m\leq\frac{8(1+\log N)}N\longrightarrow0.
\]
Thus
\[
\boxed{\displaystyle\lim_{n\to\infty}S_n=2}.
\]
\end{solution}

\begin{takeaways}
\item A floor recurrence can encode a breadth-first traversal of the Calkin--Wilf tree.
\item The matrices $M_0$ and $M_1$ encode the two Euclidean subtraction branches and give unique binary addresses to positive rationals.
\item Tree telescoping replaces cancellation along a chain by cancellation from each parent against its two children.
\item Positive series may be re-indexed by the geometry of coprime lattice points, turning the limit into a boundary estimate.
\end{takeaways}